\section{The two estimates in the smoothing at a cubic branch}\label{app:joined-estimates}
We use the parameters $h,(q_i)$, the nodal derivative $J$, and its
remaining real adjoint $\Lambda$ from \S\ref{p14-sec:joins}.
The different orders of the matching and adjoint equations determine
the degree at a join. Here we prove \eqref{p14-eq:fullestimate}
and \eqref{p14-eq:rowestimate} for the stabilized equation.

Choose thick stabilizing forms representing $P_I$ away from the
moving disks and fix a real normalization of the height, equal to one
on $s$. Solving the stabilized equation and taking its forcing
coordinate gives a continuous column $F(h,(q_i))$ with values in
$P_I$. The model \eqref{p14-eq:cubicpolynomial} permits computation in
a fixed vector bundle over $C$. Values at the three roots identify
its sections with sections of the pullback away from the branches.
The conjugates $\bar z,\bar z^2$ are replaced by their polynomial
interpolants of degree at most two. At the central fiber this imposes
the ordinary matching of the two normalizations at the node. For
$2<p_0<4$ these matrices are continuous as $L^{p_0}$ multipliers
in the parameters, by the divided-difference bound
$O(1+|w+p^2/3|^{-1/2})$. Bounded sections on the source extend over
the branches. Interpolation identifies the central stabilized kernel
by even and odd decomposition on the double sheet and matching
to the third sheet. At the smooth branch factors the same argument
applies, and the indices agree by the complex-linear finite projection formula.
Forms on thick unramified disks span both cokernels by adjoint
continuation and the matching sequence. Choose a minimal common
supplement. Both stabilized problems have a line of solutions nearby
by continuity of the pullback stabilization \cite[Appendix~B]{bai}.
Evaluation identifies the lines and their forcing coordinates.
The details of this comparison for the algebra \eqref{p14-eq:cubicpolynomial} are in \S\ref{p10-sec:nodes}.

Set $Q_i=q_i\log(1/|q_i|)$, with $Q_i=0$ when $q_i=0$.
First, in the entire supplement one has
\begin{equation}\label{p14-eq:fullestimate}
 F=Jh+2\iota\big((\overline Q_i B_i v_i)_{i\in I}\big)
                     +o(|h|+\max_i|Q_i|).
\end{equation}
Here is the coefficient of the logarithm. On a single disk put
$\ell=\max(|p|,|q|^{1/3})$. Outside $|w|\le A\ell^2$, take
$r^2=w$ and expand the roots in $p/r,q/r^3$. The quadratic
coefficient of the interpolant of $\bar z$ has linear $q$-term
\[
  -\frac32\left(\frac{\bar q}{|w|^2}
                   -\frac{\bar r}{r}\frac{q}{w^2}\right).
\]
Only its first summand has nonzero angular mean. For a fixed
H\"older test the integrated logarithmic term is
$-2\pi\bar Q\delta_0$ in the $z^2$ component. Indeed the inner
integral is $O(\ell^3)$, the outer quadratic error is
$O(|p|^2+|pq|/\ell+|q|^2/\ell^3)$, and if $|q|<|p|^3$ the
logarithmic terms are already $o(|p|)$. All other entries have
integrable $q$-derivatives. The full matrix difference has $L^1$
norm $O(|p|+|Q|)$. Continuous convergence of the normalized
stabilized solution in $W^{1,p_0}\subset C^0$ therefore permits
using its central value in this pairing to first order. The section
$\chi z^2 S(w)a/w$, with $S(0)=1$ a local real frame of solutions,
has ordinary forcing of class $\iota(a)$ and point forcing
$\pi a z^2\delta_0$. Their sum pairs to zero with a central
adjoint. Apply this to the generator with derivative $v_i$.
The disks are disjoint, which proves \eqref{p14-eq:fullestimate}.

The missing row requires a smaller remainder. Put
$R_0=\max_i|q_i|$. On each cone $|h|\le C R_0\log(1/R_0)$ on
which all $|q_i|$ are comparable and positive, one has
\begin{equation}\label{p14-eq:rowestimate}
                       \Lambda F=\sum_{i\in I} H_i(q_i)+o(R_0).
\end{equation}
We give the estimate on one disk. The half-trace of the forcing
$a_0+a_1z+a_2z^2$ is $a_0+w a_2$. Write $\varepsilon=|q|^{1/3}$,
$R=|w|^{1/2}$. On $|w|\le A\varepsilon^2$ the roots are
$O(\varepsilon)$ and their diameter is bounded below by
$c\varepsilon$: their mean is $O(p)=o(\varepsilon)$ and their
product is $-q$. The constant and quadratic coefficients of the
interpolant of $\bar z$ are $O(\varepsilon)$ and $O(1/\varepsilon)$.
The weighted inner integral is therefore $O(\varepsilon^5)$.
The central odd height is $z u_1(w)$ with $u_1(0)=v_i$ and
$u_1(w)-v_i=O(|w|)$, by oddness and the pulled-back equation.
On the exterior the weighted quadratic error from the three roots is
$O(R(|p|/R+|q|/R^3)^2)$. Integration against $O(R^3\,dR)$ gives
$O(|p|^2+|pq|+|q|^2/\varepsilon)=o(|q|)$. The linear terms are bounded
by $C(|p|+|q|/|w|)$; restoring their inner integrals costs $o(|q|)$.
Smooth center and time changes of the coefficients and collars have
Taylor remainder \(O(|h|^2)\). Their products with the varying
interpolation matrix have integral \(O(|h|(|h|+R_0\log(1/R_0)))\)
by the full matrix bound. Both are \(o(R_0)\) on the comparable cone. Other simple branch centers may be fixed by
translation on their disjoint source disks. Without translation the
angular multiplier on such a disk has first and second derivative
$O(|w|^{-1})$ and $O(|w|^{-2})$, giving integrated Taylor error
$O(|h|^2(1+|\log|h||))$.

For the varying generator, the worst unweighted matrix change in
$L^{p_0}$ on a joining disk is bounded by
\[
 C\big(\varepsilon^{4/p_0-1}
       +|p|\varepsilon^{4/p_0-2}\big).
\]
Use the root diameter on the inner disk and the outer bound
$CR^{-1}(|p|/R+|q|/R^3)$ elsewhere. Thus the bounded normalized
stabilized inverse gives $O(R_0^\delta)$ convergence of the generator
in $W^{1,p_0}$ on the comparable cone, for some $\delta>0$.
The product with the $L^1$ multiplier variation
$O(|h|+R_0\log(1/R_0))$ is $o(R_0)$ in the adjoint pairing; the
principal-part variation on the thick collars satisfies the same
bound. This justifies differentiation of the frozen height in
the half-trace pairing. At fixed regular target point the two
sextic roots move by $-q/(2w)$ to first order. Their averaged
variation has pole $-qv_i/(2w)$ at $y_i$ on the quotient cubic.
Away from that point the differentiated forcing is the cubic
operator applied to this sheet variation, including its collar cutoff.
Green's formula against the adjoint gives minus $\pi$ times
the real pairing with the residue, since
$\partial_{\bar w}(1/w)=\pi\delta_0$. The bounded remainder has
zero limiting boundary pairing. The $h$-derivative pairs to zero
by $\Lambda J=0$. Adding the integrals proves \eqref{p14-eq:rowestimate}.
